%%=============================================================================
%% Conclusie
%%=============================================================================

\chapter{Conclusie}%
\label{ch:conclusie}

Deze bachelorproef onderzocht hoe een gefaseerde migratie naar een open-source ERP-systeem zoals Odoo servicegerichte Belgische KMO’s kan ondersteunen bij het verbeteren van administratieve processen en het voorbereiden op verplichte elektronische facturatie via Peppol.

De centrale onderzoeksvraag luidde als volgt:

\begin{quote}
Hoe kan een stapsgewijze aanpak voor de migratie naar open-source ERP (Odoo) KMO’s zoals Hanuman BV helpen hun doorlooptijd en foutgraad te verlagen en tegelijk B2B-e-invoicing-compliance (Peppol) te waarborgen?
\end{quote}

Om deze onderzoeksvraag te beantwoorden werd eerst een analyse uitgevoerd van de bestaande werking van Hanuman BV. Daarbij werd vastgesteld dat de onderneming sterk afhankelijk was van Excelbestanden, e-mail en losse administratieve documenten. Deze gefragmenteerde werkwijze leidde tot dubbele gegevensinvoer, beperkte traceerbaarheid, manuele voorraadopvolging en foutgevoelige facturatieprocessen.

Op basis van de AS-IS analyse werden functionele en niet-functionele vereisten opgesteld voor een geïntegreerde ERP-oplossing. Vervolgens werd een TO-BE situatie ontworpen waarin verkoop, voorraadbeheer, serviceprocessen en facturatie centraal beheerd worden binnen Odoo.

Daarnaast werd een gefaseerd migratiekader uitgewerkt waarbij de implementatie stapsgewijs verloopt. Deze aanpak beperkt implementatierisico’s en laat toe om processen gecontroleerd te stabiliseren alvorens bijkomende functionaliteiten worden toegevoegd.

Uit de evaluatie blijkt dat de voorgestelde ERP-omgeving een duidelijke meerwaarde kan bieden voor Hanuman BV. Vooral de centralisatie van stamdata, de automatische koppeling tussen documenten en de integratie van voorraad- en serviceprocessen zorgen voor belangrijke operationele verbeteringen.

Binnen de uitgewerkte scenario’s werd vastgesteld dat:
\begin{itemize}
    \item het aantal manuele handelingen aanzienlijk vermindert;
    \item documenttraceerbaarheid sterk verbetert;
    \item voorraadbeheer actueler en consistenter wordt;
    \item facturatie sneller en minder foutgevoelig verloopt;
    \item de onderneming technisch beter voorbereid is op Peppol-compliance.
\end{itemize}

De evaluatie toont daarnaast aan dat een gefaseerde implementatie bijzonder geschikt is voor kleine ondernemingen met beperkte IT-capaciteit. Door processen stapsgewijs te digitaliseren blijft de implementatie beheersbaar en kunnen gebruikers geleidelijk vertrouwd raken met de nieuwe werkwijze.

Tegelijkertijd bevestigt dit onderzoek dat een ERP-implementatie niet uitsluitend een technische oefening is. Datakwaliteit, centrale stamdata, gebruikersacceptatie en duidelijke processtructuren vormen cruciale voorwaarden voor een succesvolle implementatie.

Hoewel de resultaten van deze bachelorproef positief zijn, kent het onderzoek ook enkele beperkingen. De evaluatie werd uitgevoerd binnen een gecontroleerde Proof of Concept-omgeving en niet binnen een volledige productieomgeving. Daarnaast waren historische logdata en langdurige gebruikersmetingen beperkt beschikbaar, waardoor bepaalde KPI’s indicatief geëvalueerd werden op basis van observaties en scenario-gebaseerde simulaties.

Verder onderzoek zou zich kunnen richten op:
\begin{itemize}
    \item langdurige productievalidatie van de implementatie;
    \item gebruikersacceptatie op langere termijn;
    \item integratie met externe boekhoud- of betaalsystemen;
    \item geavanceerde rapportering en managementdashboards;
    \item volledige productie-integratie van Peppol e-invoicing.
\end{itemize}

Samenvattend kan geconcludeerd worden dat een gefaseerde implementatie van Odoo een haalbare en zinvolle strategie vormt voor servicegerichte Belgische KMO’s zoals Hanuman BV. De combinatie van procesintegratie, centraal databeheer en voorbereiding op elektronische facturatie biedt belangrijke opportuniteiten voor operationele efficiëntie en toekomstige digitale compliance.